
\chapter{Aplicativo}
\label{chap:aplicativo}

% Apresentar o aplicativo do DoseCerta, seu objetivo
% e sua função na interação entre o usuário e o sistema.

\section{Requisitos e Objetivos da Interface}
\label{sec:requisitos-aplicativo}

% Descrever os requisitos da interface, considerando o público-alvo,
% as tarefas que o usuário precisa realizar e a facilidade de uso.

\section{Tecnologias e Ambiente de Desenvolvimento}
\label{sec:tecnologias-aplicativo}

% Informar as tecnologias, linguagens, frameworks e ferramentas
% efetivamente utilizados no desenvolvimento do aplicativo.
% Justificar as escolhas, quando pertinente.

\section{Organização e Navegação}
\label{sec:organizacao-aplicativo}

% Descrever a estrutura de navegação e a organização das telas.
% Inserir diagrama de navegação, se disponível.

\section{Telas e Funcionalidades}
\label{sec:telas-funcionalidades}

% Apresentar as telas implementadas e explicar sua finalidade.
% Inserir capturas de tela e descrever as principais interações.

% Exemplo de inserção de figura:
% \begin{figure}[htbp]
%     \centering
%     \includegraphics[width=0.8\textwidth]{Figuras/nome-da-tela}
%     \caption{Tela do aplicativo.}
%     \label{fig:tela-aplicativo}
% \end{figure}

\section{Configuração dos Medicamentos e Horários}
\label{sec:configuracao-medicamentos}

% Descrever como o usuário configura medicamentos, doses
% e horários, caso essas funcionalidades estejam disponíveis.
% Explicar validações e confirmação das configurações.

\section{Alertas e Acompanhamento}
\label{sec:alertas-acompanhamento}

% Descrever como o aplicativo apresenta lembretes,
% confirmações ou informações sobre a dispensação,
% conforme as funções efetivamente implementadas.

\section{Comunicação com o Backend}
\label{sec:comunicacao-aplicativo-backend}

% Explicar como o aplicativo troca dados com o backend.
% Informar endpoints, formatos de dados e tratamento de falhas
% relevantes para a interface, sem repetir todo o Capítulo 10.

\section{Acessibilidade e Usabilidade}
\label{sec:acessibilidade-usabilidade}

% Discutir decisões de interface pertinentes ao público-alvo,
% como legibilidade, tamanho de elementos, clareza dos comandos
% e facilidade de navegação.
% Distinguir decisões de projeto de avaliações efetivamente realizadas.

\section{Síntese do Aplicativo}
\label{sec:sintese-aplicativo}

% Resumir as funções disponíveis, sua integração com o DoseCerta
% e as limitações atuais do aplicativo.
